\chapter{نظام انتخاباتی و مشارکت دموکراتیک}
\label{app:elections}

طراحی سیستم انتخاباتی تعیین‌کننده‌ی تکثرگرایی و ثبات در آینده‌ی ایران است.

\section{نظام مختلط تناسبی (MMP)}

طرح مهستان برای مجلس دائمی (و با تغییراتی برای مجلس مهستان) سیستم MMP (مشابه آلمان) را پیشنهاد می‌دهد.

\begin{modernbox}{ساختار رأی‌دهی MMP}
هر شهروند دو رأی می‌دهد:
\begin{itemize}
    \item \textbf{رأی اول (حوزه‌ای):} انتخاب یک فرد برای نمایندگی محلی (سیستم اکثریتی). این امر ارتباط مستقیم نماینده و مردم را حفظ می‌کند.
    \item \textbf{رأی دوم (فهرستی):} انتخاب یک جریان سیاسی (سیستم نسبی). این رأی سهم کلی هر جریان در مجلس را تعیین می‌کند.
\end{itemize}
کرسی‌های فهرستی به گونه‌ای توزیع می‌شوند که نسبت نهایی صندلی‌ها در مجلس با درصد آرای فهرستی جریان‌ها برابری کند.
\end{modernbox}

\section{توزیع کرسی‌ها در مجلس مهستان (۳۰۰ کرسی)}

\begin{center}
\begin{tikzpicture}
    \pie[rotate=0, radius=3, text=inside, color={LightBlue, SoftGreen, SoftPurple, SoftGold}]
    {150/کرسی‌های فهرستی ملی, 130/کرسی‌های حوزه‌ای استانی, 20/کرسی‌های تضمینی اقلیت‌ها}
\end{tikzpicture}
\end{center}

\section{حقوق انتخاباتی خاص}

\begin{itemize}
    \item \textbf{سهمیه‌ی جنسیتی:} حداقل ۳۰٪ از کاندیداهای هر فهرست باید از زنان باشند (با سیستم زیپ: یک در میان).
    \item \textbf{کرسی‌های اقلیت‌ها:} ۲۰ کرسی تضمینی برای اقلیت‌های قومی و مذهبی که با سیستم اکثریتی در مناطق خود انتخاب می‌شوند.
    \item \textbf{رأی دیاسپورا:} ایجاد حوزه‌های انتخابیه مجازی یا کنسولی برای میلیون‌ها ایرانی خارج از کشور.
\end{itemize}

\section{ملاحظات فنی رأی‌گیری دیجیتال}

\begin{modernbox}{بستر دموکراسی مستقیم}
رأی‌گیری دیجیتال در طرح مهستان \textbf{مکمل} است، نه جایگزین رأی‌گیری حضوری.
\begin{itemize}
    \item استفاده از بلاک‌چین برای شفافیت و حسابرسی‌پذیری.
    \item احراز هویت دوعاملی (کد ملی + بیومتریک اختیاری).
    \item امکان نظرسنجی مداوم و طومارهای الکترونیکی برای عزل نماینده.
\end{itemize}
\end{modernbox}

\section{مقایسه سیستم‌های انتخاباتی برای ایران}

\begin{center}
\captionof{table}{ارزیابی سیستم‌های انتخاباتی}
\begin{tabular}{R{3cm} C{3cm} C{3cm} C{3cm}}
\hline
\textbf{معیار} & \textbf{اکثریتی (FPTP)} & \textbf{نسبی فهرستی (PR)} & \textbf{مختلط (MMP)} \\
\hline
%\endfirsthead
%\hline
%\textbf{معیار} & \textbf{اکثریتی (FPTP)} & \textbf{نسبی فهرستی (PR)} & \textbf{مختلط (MMP)} \\
%\hline
%\endhead
%\hline
%\endfoot
سادگی & بسیار بالا & بالا & متوسط \\
عدالت در نمایندگی & پایین & بسیار بالا & بسیار بالا \\
ارتباط نماینده-حوزه & بسیار قوی & ضعیف & قوی \\
ثبات دولتی & بالا & متوسط & بالا \\
مناسب برای ایران & خیر & بله & \textbf{عالی} \\
\hline
\end{tabular}
\end{center}
